\begin{abstract}
Evaluations of attribute--object binding in vision--language models can be invalidated in ways
that are easy to miss: the model may receive the binding through its input features, a metric may
be solvable by a binding-blind content signal, or a single-modality control may look harmless only
because of how the metric is constructed. We describe an evaluation setting for compositional
binding in which a scene-graph generator produces minimal-change groups of two images and two
captions, every group is validated by an oracle over the scene graph, and the model sees only
rendered pixels and raw caption strings. Splits keep each edit orbit (all scenes with the same
multisets of shapes, colours and materials, which differ only in how these are assigned to objects
and positions) inside a single split, and the endpoint is restricted to binding-necessary groups, whose two scenes lie in
the same orbit. We use this setting to examine a small trainable head on frozen CLIP ViT-B/32
tokens trained with a hard-negative contrastive loss. On a development panel of \SPtrainGroups{}
training and \SPdevGroups{} development groups, the head largely fits its training groups
(binding-necessary training group score \SPtrainBNGroupMin--\SPtrainBNGroupMax), but on the
\SPdevBNn{} binding-necessary development groups its group score ranges from \SPdevBNGroupMin{} to
\SPdevBNGroupMax{} over three optimization repeats (the chance level for independent continuous
scores is $1/6$), and most of its full-panel score is reproduced by a binding-blind content channel.
We also show that single-modality scorers are structurally zero on $2{\times}2$ group metrics and
exactly $0.5$ on balanced pair AUC, so they cannot certify the absence of shortcuts.
\ifhavepaired
In a paired comparison on a larger panel (\SSpanelTrainN{} training and \SSpanelDevN{} development
groups, three optimization repeats), the hard-negative loss was compared with the same loss plus a
cross-group edit-consistency objective. Neither arm met a pre-fixed training-fit criterion, and
training was unstable in both. The differences in favour of the edit objective appear only in runs
where the baseline ended with a high final loss, so we record the comparison as unresolved and put
that branch on hold.

\else
A paired comparison of the hard-negative loss with an added edit-consistency objective has not been
run in this draft.
\fi
All results concern synthetic renders, a frozen encoder and development data; the held-out test
splits remain unopened.
\end{abstract}
